\section{Exercises}


\begin{frame}
\frametitle{Strings}
Given the string 'hello' give an index command that returns 'e'.
\end{frame}

\begin{frame}[fragile]
\frametitle{Strings}
Reverse the string 'hello' using slicing
\pause
 \begin{lstlisting}[language=Python]
    s  = 'hello'
    s[::-1]
\end{lstlisting}
\end{frame}

\begin{frame}[fragile]
\frametitle{Strings}
Given the string hello, give two methods of producing the letter 'o' using indexing.
\pause
 \begin{lstlisting}[language=Python]
    s = 'hello'
    s[-1]
    s[4]
\end{lstlisting}
\end{frame}

\begin{frame}[fragile]
\frametitle{Lists}
Build this list [0,0,0] two separate ways.
\pause
 \begin{lstlisting}[language=Python]
    [0]*3
    [0,0,0]
\end{lstlisting}
\end{frame}

\begin{frame}[fragile]
\frametitle{Lists}
Reassign 'hello' in this nested list to say 'goodbye' instead: 
 \begin{lstlisting}[language=Python]
    list3 = [1,2,[3,4,'hello']]
\end{lstlisting}
\pause
 \begin{lstlisting}[language=Python]
    list3[2][2] = 'goodbye'
\end{lstlisting}
\end{frame}

\begin{frame}[fragile]
\frametitle{Lists}
Sort the list below: 
 \begin{lstlisting}[language=Python]
    list4 = [5,3,4,6,1]
\end{lstlisting}
\pause
 \begin{lstlisting}[language=Python]
    list4.sort()
\end{lstlisting}
\end{frame}

\begin{frame}[fragile]
\frametitle{Dictionaries}
Using keys and indexing, grab the 'hello' from the following dictionaries:
 \begin{lstlisting}[language=Python]
    d = {'simple_key':'hello'}
    # Grab 'hello'
    d = {'k1':{'k2':'hello'}}
    d = {'k1':[{'nest_key':['this is deep',['hello']]}]}
    d = {'k1':[1,2,{'k2':['this is tricky',{'tough':[1,2,['hello']]}]}]}
\end{lstlisting}
\pause
 \begin{lstlisting}[language=Python]
    d['simple_key']
    d['k1']['k2']
    d['k1'][0]['nest_key'][1][0]
    d['k1'][2]['k2'][1]['tough'][2][0]
\end{lstlisting}
\end{frame}

\begin{frame}
\frametitle{Dictionaries}
Can you sort a dictionary? Why or why not? 
\pause 
\begin{itemize}
        \item \textbf{Mappings vs. Sequences:}
        \begin{itemize}
            \item Dictionaries are mappings, not sequences.
            \item Designed to map keys to values, rather than maintain a specific order.
        \end{itemize}
        \item \textbf{Hash Table Implementation:}
        \begin{itemize}
            \item Implemented using hash tables.
            \item Provides average-case O(1) time complexity for lookups, insertions, and deletions.
        \end{itemize}
        \item \textbf{Order Preservation (Python 3.7+):}
        \begin{itemize}
            \item Starting with Python 3.7, dictionaries maintain insertion order.
            \item Guaranteed by the language specification in Python 3.8+.
        \end{itemize}
        \item \textbf{Sorting Representations:}
        \begin{itemize}
            \item Dictionaries are not inherently sorted.
            \item Create sorted representations using sorted lists or \texttt{OrderedDict}.
        \end{itemize}
\end{itemize}
\end{frame}

\begin{frame}
\frametitle{Tuples}
What is the major difference between tuples and lists?
\begin{itemize}
    \item \textbf{Mutability:}
    \begin{itemize}
        \item Lists are mutable, while tuples are immutable.
        \item Lists can be modified after creation, while tuples cannot.
        \item Tuples are hashable, while lists are not.
    \end{itemize}
\end{itemize}
\end{frame}

\begin{frame}
\frametitle{Sets}
What is unique about a set?
    \begin{itemize}
        \item \textbf{Unordered Collection:}
        \begin{itemize}
            \item Elements do not have a specific order.
            \item No indexing or slicing.
        \end{itemize}
        \item \textbf{Unique Elements:}
        \begin{itemize}
            \item All elements are unique.
            \item Duplicate elements are ignored.
        \end{itemize}
        \item \textbf{Mutable:}
        \begin{itemize}
            \item Elements can be added or removed.
        \end{itemize}
    \end{itemize}
\end{frame}

\begin{frame}
\frametitle{Sets}
Use a set to find the unique values of the list below:
\end{frame}



\begin{frame}[fragile]
\frametitle{Booleans}
What will be the resulting Boolean of the following pieces of code 
 \begin{lstlisting}[language=Python]
    2 > 3
    3 <= 2
    3 == 2.0
    4**0.5 != 2
    # two nested lists
l_one = [1,2,[3,4]]
l_two = [1,2,{'k1':4}]
    # True or False?
l_one[2][0] >= l_two[2]['k1']
    \end{lstlisting}
\end{frame}

\begin{frame}[fragile]
\frametitle{Expression Evaluation}
Please tell the order in which the Python interpreter would evaluate them
 \begin{lstlisting}[language=Python]
x = 2
y = 3
square(sub((1+y),x)) 
\end{lstlisting}
\end{frame} 


\begin{frame}{Answer}
\begin{enumerate}
    \item look up the variable to get the square function Object
    \item look up the variable to get the sub function Object
    \item look up the variable to get the y value = 3
    \item add 1 and 3 to get 4
    \item look up the variable to get the x value = 2
    \item run the sub function passing 4 and 2 returning 2
    \item run the square function passing 2 returning 4
\end{enumerate}
\end{frame}

\begin{frame}[fragile]
\frametitle{Expression Evaluation}
Please tell the order in which the Python interpreter would evaluate them
 \begin{lstlisting}[language=Python]
x = 2
y = 3
square(x + sub(square(y), 2 *x))
\end{lstlisting}
\end{frame} 




\begin{frame}[fragile]{Naming Your Robot}
    \begin{itemize}
        \item Use \texttt{input()} to read the robot's name from the user.
        \item Assign the name to a variable called \texttt{name}.
        \item Assign an identifier to the robot, e.g., \texttt{identifier = 1000}.
        \item Print a message from the robot using \texttt{print()}.
    \end{itemize}
    \begin{lstlisting}[language=Python]
>>> name = input("Enter the robot's name: ")
>>> identifier = 1000
>>> print(f"Hello. My name is {name}. My ID is {identifier}.")
    \end{lstlisting}
\end{frame}

\begin{frame}[fragile]{Converting 2D Index to 1D Index}
    \begin{itemize}
        \item User provides the number of rows \(R\), number of columns \(C\), and a 2D index \((r, c)\)
        \item Convert the 2D index \((r, c)\) to a 1D index \(s\) using row-major ordering.
    \end{itemize}
    \begin{lstlisting}[language=Python]
# Get user input
R = int(input("Enter the number of rows: "))
C = int(input("Enter the number of columns: "))
r = int(input("Enter the row index: "))
c = int(input("Enter the column index: "))
s = r*R + c
# Print the result
print(f"The 1D index is: {s}")
    \end{lstlisting}
\end{frame}

\begin{frame}[fragile]{Infinite Monkey Theorem}
 \begin{lstlisting}[language=Python]
import string 
import random 
char_dictionary = string.ascii_lowercase+' '
target = 'hello world'
def genrate_random_string(length):
    random_string = ''
    for i in length:
        random_string += random.choice(char_dictionary)
    return random_string
\end{lstlisting}
\end{frame}

\begin{frame}[fragile]{Infinite Monkey Theorem}
 \begin{lstlisting}[language=Python]
def infinite_monkey(target):
    target_length = len(target)
    iteration = 0
    while True:
        random_string = genrate_random_string(target_length)
        iteration += 1
        print(iteration)
        if random_string == target:
            return iteration
attempts = infinite_monkey(target)
\end{lstlisting}
\end{frame}


